% GENERATED FILE. Edit canonical lessons, then run npm run generate:books.
% canonical-source-hash: fnv1a64:34e2c8ace3d8e6db
% canonical-lessons: BN-C316-ghera, BN-C316-prachir, BN-C316-porikha, BN-C316-teka, BN-C316-pichhono

\chapter{The Siege}
\label{ch:bn-the-siege}
\hlchaptermodality{316}

\begin{chapteropening}
\textbf{By the end of this chapter:} \emph{I can describe a fort under siege: walls, moats, holding out.}

Complete the last lesson of chapter 316: I can describe a fort under siege: walls, moats, holding out.
\end{chapteropening}

\section[\texorpdfstring{\bn{ঘেরা} (gherā)}{ঘেরা (gherā)}]{\bn{ঘেরা} (gherā) — to surround; to besiege}

\label{lesson:BN-C316-ghera}



\begin{quote}
Before the new one: say the Bengali for a gun, then the Bengali for a cannon.
\end{quote}

\subsection*{You'll want to know: \bn{ঘেরা}}
\textbf{\bn{ঘেরা}} — \emph{gherā} — \textquotedblleft{}to surround; to besiege\textquotedblright{}.

\begin{quote}
\bn{শত্রুর} \bn{সৈন্যরা} \bn{দুর্গটা} \bn{ঘিরে} \bn{ফেলল।}
\end{quote}

\begin{quote}
\emph{shôtrur soinyôrā durgôṭā ghire phellô.}
\end{quote}

\begin{quote}
The enemy soldiers surrounded the fort.
\end{quote}

Past (he/she …): \bn{সেই} \bn{সেনাবাহিনী} \bn{শহরটা} \bn{ঘিরল}, \emph{sei senābāhinī shôhôrṭā ghirlô}.

The stem vowel rises in the past: \bn{ঘিরল}, surrounded.

\begin{cousinweb}
\textbf{\bn{ঘেরা}} (\emph{gherā}) is filed in CDIAL 4474 under a Middle Indo-Aryan \emph{ghir-}, 'go round', and \emph{ghēr-}, 'surround', which list Bengali \emph{gherā}, 'to surround'.
\end{cousinweb}

\begin{grammarlens}[title={What you've built}]
One more word for telling what happened.
\end{grammarlens}

\subsection*{Your turn}
\begin{itemize}
\raggedright
  \item \emph{Say it:} \emph{gherā}
  \item \emph{Say it:} \emph{gherā}, once more
  \item \emph{Say it:} \emph{shôtrur soinyôrā durgôṭā ghire phellô.} — The enemy soldiers surrounded the fort.
  \item \emph{Recall:} say the Bengali for a gun, then the Bengali for a cannon, then say \emph{gherā} again
\end{itemize}

\begin{tcolorbox}[breakable,colback=teal!4,colframe=teal!35!black,title={Before you move on}]
What does \bn{ঘেরা} mean? (\textquotedblleft{}To surround; to besiege\textquotedblright{}.) Say it once more.
\end{tcolorbox}
\section[\texorpdfstring{\bn{প্রাচীর} (prāchīr)}{প্রাচীর (prāchīr)}]{\bn{প্রাচীর} (prāchīr) — a rampart; a city wall}

\label{lesson:BN-C316-prachir}



\begin{quote}
Before the new one: say the Bengali for a cannon, then the Bengali for to surround.
\end{quote}

\subsection*{You'll want to know: \bn{প্রাচীর}}
\textbf{\bn{প্রাচীর}} — \emph{prāchīr} — \textquotedblleft{}a rampart; a city wall\textquotedblright{}.

\begin{quote}
\bn{শহরের} \bn{চারপাশে} \bn{একটা} \bn{উঁচু} \bn{প্রাচীর} \bn{ছিল।}
\end{quote}

\begin{quote}
\emph{shôhôrer chārpāshe ekṭā ũchu prāchīr chhilô.}
\end{quote}

\begin{quote}
A high wall ran all round the city.
\end{quote}

\begin{cousinweb}
\textbf{\bn{প্রাচীর}} (\emph{prāchīr}) is Sanskrit \emph{prācīra}, 'an enclosure, a fence, a wall' (Monier-Williams); a house wall is \textbf{\bn{দেয়াল}}.
\end{cousinweb}

\begin{grammarlens}[title={What you've built}]
One more word for telling what happened.
\end{grammarlens}

\subsection*{Your turn}
\begin{itemize}
\raggedright
  \item \emph{Say it:} \emph{prāchīr}
  \item \emph{Say it:} \emph{prāchīr}, once more
  \item \emph{Say it:} \emph{shôhôrer chārpāshe ekṭā ũchu prāchīr chhilô.} — A high wall ran all round the city.
  \item \emph{Recall:} say the Bengali for a cannon, then the Bengali for to surround, then say \emph{prāchīr} again
\end{itemize}

\begin{tcolorbox}[breakable,colback=teal!4,colframe=teal!35!black,title={Before you move on}]
What does \bn{প্রাচীর} mean? (\textquotedblleft{}A rampart; a city wall\textquotedblright{}.) Say it once more.
\end{tcolorbox}
\section[\texorpdfstring{\bn{পরিখা} (pôrikhā)}{পরিখা (pôrikhā)}]{\bn{পরিখা} (pôrikhā) — a moat; a trench}

\label{lesson:BN-C316-porikha}



\begin{quote}
Before the new one: say the Bengali for to surround, then the Bengali for a rampart.
\end{quote}

\subsection*{You'll want to know: \bn{পরিখা}}
\textbf{\bn{পরিখা}} — \emph{pôrikhā} — \textquotedblleft{}a moat; a trench\textquotedblright{}.

\begin{quote}
\bn{প্রাচীরের} \bn{বাইরে} \bn{জলে} \bn{ভরা} \bn{একটা} \bn{পরিখা} \bn{ছিল।}
\end{quote}

\begin{quote}
\emph{prāchīrer bāire jôle bhôrā ekṭā pôrikhā chhilô.}
\end{quote}

\begin{quote}
Outside the wall there was a moat full of water.
\end{quote}

\begin{cousinweb}
\textbf{\bn{পরিখা}} (\emph{pôrikhā}) is Sanskrit \emph{parikhā}, 'a moat, a ditch round a fort' (Monier-Williams): \emph{pari-}, 'around', and \emph{√khan}, 'to dig'.
\end{cousinweb}

\begin{grammarlens}[title={What you've built}]
One more word for telling what happened.
\end{grammarlens}

\subsection*{Your turn}
\begin{itemize}
\raggedright
  \item \emph{Say it:} \emph{pôrikhā}
  \item \emph{Say it:} \emph{pôrikhā}, once more
  \item \emph{Say it:} \emph{prāchīrer bāire jôle bhôrā ekṭā pôrikhā chhilô.} — Outside the wall there was a moat full of water.
  \item \emph{Recall:} say the Bengali for to surround, then the Bengali for a rampart, then say \emph{pôrikhā} again
\end{itemize}

\begin{tcolorbox}[breakable,colback=teal!4,colframe=teal!35!black,title={Before you move on}]
What does \bn{পরিখা} mean? (\textquotedblleft{}A moat; a trench\textquotedblright{}.) Say it once more.
\end{tcolorbox}
\section[\texorpdfstring{\bn{টেকা} (ṭekā)}{টেকা (ṭekā)}]{\bn{টেকা} (ṭekā) — to hold out; to last}

\label{lesson:BN-C316-teka}



\begin{quote}
Before the new one: say the Bengali for a rampart, then the Bengali for a moat.
\end{quote}

\subsection*{You'll want to know: \bn{টেকা}}
\textbf{\bn{টেকা}} — \emph{ṭekā} — \textquotedblleft{}to hold out; to last\textquotedblright{}.

\begin{quote}
\bn{দুর্গটা} \bn{তিন} \bn{মাস} \bn{টিকল।}
\end{quote}

\begin{quote}
\emph{durgôṭā tin mās ṭiklô.}
\end{quote}

\begin{quote}
The fort held out for three months.
\end{quote}

Past (he/she …): \bn{সেই} \bn{দুর্গ} \bn{টিকল}, \emph{sei durgô ṭiklô}.

The stem vowel rises in the past: \bn{টিকল}, held out.

\begin{cousinweb}
\textbf{\bn{টেকা}} (\emph{ṭekā}) is filed in CDIAL 5420 under a Middle Indo-Aryan \emph{ṭakk-}, 'remain, stop', which lists Bengali \emph{ṭẽkā} and \emph{ṭikā}, 'to remain'.
\end{cousinweb}

\begin{grammarlens}[title={What you've built}]
One more word for telling what happened.
\end{grammarlens}

\subsection*{Your turn}
\begin{itemize}
\raggedright
  \item \emph{Say it:} \emph{ṭekā}
  \item \emph{Say it:} \emph{ṭekā}, once more
  \item \emph{Say it:} \emph{durgôṭā tin mās ṭiklô.} — The fort held out for three months.
  \item \emph{Recall:} say the Bengali for a rampart, then the Bengali for a moat, then say \emph{ṭekā} again
\end{itemize}

\begin{tcolorbox}[breakable,colback=teal!4,colframe=teal!35!black,title={Before you move on}]
What does \bn{টেকা} mean? (\textquotedblleft{}To hold out; to last\textquotedblright{}.) Say it once more.
\end{tcolorbox}
\section[\texorpdfstring{\bn{পিছোনো} (pichhono)}{পিছোনো (pichhono)}]{\bn{পিছোনো} (pichhono) — to fall back; to retreat}

\label{lesson:BN-C316-pichhono}



\begin{quote}
Before the new one: say the Bengali for a moat, then the Bengali for to hold out.
\end{quote}

\subsection*{You'll want to know: \bn{পিছোনো}}
\textbf{\bn{পিছোনো}} — \emph{pichhono} — \textquotedblleft{}to fall back; to retreat\textquotedblright{}.

\begin{quote}
\bn{খাবার} \bn{শেষ} \bn{হলে} \bn{শত্রুরা} \bn{পিছোল।}
\end{quote}

\begin{quote}
\emph{khābār shesh hôle shôtrurā pichholô.}
\end{quote}

\begin{quote}
When their food ran out, the enemy fell back.
\end{quote}

Past (he/she …): \bn{সেই} \bn{সেনাবাহিনী} \bn{পিছোল}, \emph{sei senābāhinī pichholô}.

\begin{cousinweb}
\textbf{\bn{পিছোনো}} (\emph{pichhono}) is built on \bn{পিছে}, behind — CDIAL 7990, a Middle Indo-Aryan \emph{paśca}, 'the hinder part', which lists Bengali \emph{pichā}, 'the rear' — as \textbf{\bn{পিছনে}} is: to move to the back.
\end{cousinweb}

\begin{grammarlens}[title={What you've built}]
One more word for telling what happened.
\end{grammarlens}

\subsection*{Your turn}
\begin{itemize}
\raggedright
  \item \emph{Say it:} \emph{pichhono}
  \item \emph{Say it:} \emph{pichhono}, once more
  \item \emph{Say it:} \emph{khābār shesh hôle shôtrurā pichholô.} — When their food ran out, the enemy fell back.
  \item \emph{Recall:} say the Bengali for a moat, then the Bengali for to hold out, then say \emph{pichhono} again
\end{itemize}

\begin{tcolorbox}[breakable,colback=teal!4,colframe=teal!35!black,title={Before you move on}]
What does \bn{পিছোনো} mean? (\textquotedblleft{}To fall back; to retreat\textquotedblright{}.) Say it once more.
\end{tcolorbox}
